\section{Q\&As: (Questions in comments) }
%Kenya's blackout due to cascading failure without a good recovery scheme
For example, a cascading failure happened in Kenya \cite{apkenya_dec23} that took 14 hours to recover and affected 50 million people \cite{apkenya_dec23}. The triggering event was an error message sent to the Lake Turkana wind power plant causing it to go offline \cite{apkenya_dec23}, losing 270 MW in power generation, which is 15\% of the local demand. This triggering event led to an imbalance of power and shut down other main generators, causing a large blackout area in Kenya, so they had to import power from a nearby country to begin the recovery process. However, after they set up the connection with Uganda's power grid, there was no power coming from Uganda \cite{apkenya_aug23}. At the same time, the failure continued to propagate throughout the country, resulting in a nationwide blackout. This example is one of many that highlight the importance of the careful selection of recovery actions and the serious consequences of cascading failures, which pose serious challenges to the stability, sustainability, and survivability of complex networked systems \cite{xiaohui_review_2016}. 

%SSS - Comparison to the recent blackout in Spain would be a good idea
%SSS - explicitly mention the Spain blackout
the 2025 Iberian Peninsula (Spain) blackout illustrates how coordinated, informed decision-making can limit damage. Operators leveraged cross-border support from France and Morocco, applied controlled load shedding, and prioritized critical infrastructure, effectively containing the disruption and speeding restoration~\cite{cetinkaya_technical_2025}.

Grid following-inverter-based resources (GFL-IBR) was the root cause of the cascading failure, as it will automatically match the renewable energy generation with the power flow in the grid. When a power error occurs in the grid, all the GFL-IBR will try to match with the error, caused the generators to shut down.
If there are more Grid forming inverter-based resources (GFM-IBR) installed in the system, the cascade could be eliminated. GFM-IBR can generate the “correct”(or say a “fixed”) quality of power for the gird. If the GFL can follow the GFM then the cascading failure could be reduced.
However, GFM-IBR are expansive and complex to make, to adjust the setting for the best quality of power. It’s a new technology, and we don’t know much about its consequences for massive installation.


To predict failure propagation paths and make mitigation designs, many mathematical models have been developed in recent years to elucidate the effect of failures and estimate the rate of recovery. 

%SSS - Introduce the concept of winning a race against failure propagation, and explain that is what we to enable here. 
When a random fault is detected in a complex networked system, it does not take long for the fault to spread to the next component via coupling. And as more and more components become failed, the rate of spread will increase faster and faster. The repair crews can start working on recovering the system while cascading failure is happening, but they have to make the right decision. If they don't know what they are doing, like the story in Kenya, they spend hours recovering an interconnection that does not help the system; in the meantime, the failure already spreads nationwide. 
Therefore, to win the race against failure propagation, we cannot just randomly repair failed components, we must target the critical ones that are essential to the system. 

Our model provided an environment for testing the recovery strategy in a complex networked system under cascading failure. Our model was designed with flexibility for various types of systems, whether they're power, water, gas, or data networks. 
They could have different interdependencies between components, but we can quantify them based on previous failure cases. These interdependencies will be reflected in failure propagation captured by our model. 

They could have different failure scenarios that could be 1, 2, or more failures happening to the system selected randomly or in a meaningful way to test the system's survivability facing different types of hazards from small, as a minor error, to large, like a natural disaster. Our model is applicable regardless of the distribution of repair time of components. 

%SSS - explain how the model is useful to others. Why should they care?
%Others include other researcher, power utilities, government decision/policy makers, etc. Anything we say here needs to be firmly backed up in the text and reiterated in the conclusions. 
%SSS - the statement immediately below describes the simulation environment, not the model. We simulate only smart grids, not other complex networks. 
Our model is valuable in various applications; it gives prevision and strategy for recovering a complex networked system. It enables deeper understanding of the race between recovery and cascading failure. It allows better design, task scheduling, and efficient use of recovery resources. 

For decision makers of the complex networked system, they can use the model to find out the most suitable way to recover the network before anything has happened, the default rules. And they adjust the policy during a cascading failure to find the best way to mitigate it before they completely lose the race.


%SSS - add a very brief explanation of what is missing from the body of knowledge that makes our work necessary. For example, "nothing exists that would help people decide how to win the race". Don't forget to check the paper recommended by the QEST reviewer - the one where failures were happening along with recovery. 

Studies in recent years are mainly focused on cascading failure analysis or recovering a complex networked system without any cascade. 
%SSS - Sanfan means that analysis of cascading failure is not enough .
%SSS - explicitly mention the paper that proposes tweaking dependencies and explain that it's costly and potentially infeasible as a solution
Analyzing the causes of cascading failure could help network managers to enhance or tweak the dependencies between components at an expensive cost, and due to the complexity of the network, cascading failures could still happen, like the grid in Spain. 
The recovery model without considering cascading failure is not useful either, as most of the complex networked systems are serving as the essential service of the society; they cannot be switched off for a post-hazard recovery or wait till the cascading failure finishes to begin recovery. The systems must be recovering while the cascading failure, and the repair crews must win the race before the failure becomes a serious problem that leads to a catastrophic failure with heavy penalties. 

Now, there are some studies that focus on recovering a complex networked system while failures occur. However, they assumed the repair time of each component is previously known, or assumed the repair time was an exponential distribution. The formal assumption is possible if the time scale is large, like if we are talking about days. People can do a lot of things in one day, and they can definitely fix whatever was broken, but the large time scale could miss lots of details in a cascading failure. It's also impossible to use on a cascading failure that recovery needs to be done in hours. The later assumption is better for cascading failure in the hourly scale, or in minutes. But it's not always true; the recovery process of a component is not always in exponential distributions. Uniform distribution and Weibull distribution are also commonly seen in the recovery time of a component, and the model should be flexible to accept those distributions. %SSS - add a reference for uniform/Weibull repair time

%SSS - refer to a specific paper for each of these statements (including the ones above)
Other than that, the recent work in recovery models has overlooked the importance of other layers of the system. For example, when recovering a smart grid, the research usually focuses on recovering the physical layer of the network, ``trying to get as much line powered as possible,'' but overlooks the importance of other layers of the system, such as the cyber layer. The cyber layer, or the control system layer, describes the cyber components of the system that adjust the flow and stabilize the quality of the service, and is an important element to consider when recovering a failing system.

%SSS - add a very brief reason for choosing a TI MDP to create what is missing from the body of knowledge
As mentioned above, some research assumes that the repair time is known, or in an exponential distribution. And we have discussed that this is not always possible or true, and we must consider other types of distribution as well. Time-inhomogeneous transitions are introduced to this model to solve the issue with different repair time distributions. 

As a Markov model, the future transition only relies on the current state and is independent of the previous states. However, the effort to repair spent on the failed components cannot be undone each time the model transitions to a new state. Also, the repair cannot be done in one time step like a miracle, so the transition for the repair completion process must be tracked by the time each crew spent on the component. Therefore, time-inhomogeneous transitions are utilized in this model to track the repair process of each component during the cascading failure. With TI, other distributions can be put into the repair time.

%SSS - need to mention that the model reflects dependencies among components, by using our previous work
This paper presents a Markov decision process (MDP) model with time-inhomogeneous transitions that reflect the recovery process of a complex networked system under cascading failure. 

The model considers dependencies between physical and cyber components as the root cause of cascading. The identification and quantification of the interdependency were present in our previous work \cite{marashi_identification_2021}, and then we created a dependency index to represent the conditional probability of failure propagation between components. When a component goes down, it could generate error messages or faults to other components that it connects to. When a component just goes down, it could bring the most impact to the system. As time continues, the impact will become less and less as the failed components would shut down or the system can bear the faults. Therefore, the fault propagation probability is going to start at the highest value when the component fails and decrease throughout the rest of the recovery process until the component gets restored and the fault propagation will no longer happen. The time-inhomogeneous transition can be used with the system state that tracks the failed moment of a component and replicates the fault propagation process. 
